% =====================================================================
%  2.1 Likninger og formler — teori og oppgaver
% =====================================================================

\del{1}{Uttrykk, likning, identitet og funksjon}

\begin{ordliste}
algebraisk uttrykk & Tall og bokstaver knyttet sammen med regneartene. Har ikke likhetstegn. & \mittsprak\\
likning & To uttrykk med likhetstegn mellom. Sann for noen verdier av $x$. & \mittsprak\\
identitet & En likning som er sann for alle tillatte verdier av $x$. & \mittsprak\\
funksjon & En regel som gir nøyaktig én verdi ut for hver verdi av $x$. & \mittsprak\\
moteksempel & Én verdi som viser at en påstand ikke alltid stemmer. & \mittsprak\\
\end{ordliste}

\begin{mustkunne}
\renewcommand{\arraystretch}{1.4}
\begin{tabular}{@{}l l l@{}}
\textbf{Begrep} & \textbf{Eksempel} & \textbf{Hva gjør vi?}\\
Uttrykk & $2x + 3$ & Forenkler. Regner ut verdien for et gitt $x$.\\
Likning & $2x + 3 = 7$ & \emph{Løser} den. Svaret er de $x$-ene som gjør den sann.\\
Identitet & $(x+1)^2 = x^2 + 2x + 1$ & \emph{Viser} den med algebra.\\
Funksjon & $f(x) = 2x + 3$ & Bruker og tegner den.\\
\end{tabular}

\medskip
\textbf{Hvorfor algebra?} Ett tall som gjør likningen usann, er et moteksempel. Da er det ikke en identitet.
Mange tall som stemmer, beviser ingenting, for neste tall kan svikte.
Når vi løser en likning og ender på $0 = 0$, er alle tall løsninger. Ender vi på en usann påstand, er det ingen løsning.
\end{mustkunne}

\begin{eksempel}[1 — er det en identitet?]
Avgjør om $(x-2)^2 = x^2 - 4$ er en identitet.

\begin{stegtabell}
Sett inn $x = 1$: & Venstre side: $(1-2)^2 = 1$. Høyre side: $1 - 4 = -3$.\\
$1 \neq -3$ & Moteksempel. Det er ikke en identitet.\\
$(x-2)^2 = x^2 - 4x + 4$ & Riktig versjon, fra kvadratsetningen.\\
\end{stegtabell}
\end{eksempel}

\begin{eksempel}[2 — løs og tolk]
Løs $2(x+3) = 2x + 6$ og $2(x+3) = 2x + 5$.

\begin{stegtabell}
$2x + 6 = 2x + 6$ & Gang ut parentesen i den første.\\
$0 = 0$ & Trekk fra $2x + 6$. Alltid sant: alle tall er løsninger. Det er en identitet.\\
$2x + 6 = 2x + 5$ & Gang ut parentesen i den andre.\\
$0 = -1$ & Trekk fra $2x + 5$. Aldri sant: ingen løsning.\\
\end{stegtabell}
\end{eksempel}

\begin{vanligfeil}
\begin{itemize}[leftmargin=1.2em,itemsep=1pt,topsep=0pt]
  \item «$0 = 0$, så $x = 0$» er feil. Likningen $2(x+3) = 2x + 6$ gir $0 = 0$, og alle tall er løsninger.
  \item «$0 = 5$, så $x = 5$» er feil. Likningen $2(x+3) = 2x + 11$ gir $0 = 5$, og har ingen løsning.
  \item «Det stemmer for to tall jeg prøvde, så det er en identitet» er feil. Noen få tall beviser ingenting.
\end{itemize}
\end{vanligfeil}

\nivaahode{1}

\oppg{1.1}{1}{\uh}
Er dette et uttrykk, en likning, en identitet eller en funksjon?
\begin{tasks}(2)
  \task $4x - 1$
  \task $4x - 1 = 11$
  \task $f(x) = 4x - 1$
  \task $3(x+4) = 3x + 12$
\end{tasks}

\oppg{1.2}{1}{\uh}
Sett inn $x = 4$. Er $x = 4$ en løsning?
\begin{tasks}(3)
  \task $3x - 5 = 7$
  \task $2x + 1 = x + 6$
  \task $x^2 = 3x + 4$
\end{tasks}

\oppg{1.3}{1}{\uh}
Du har forenklet en likning og står igjen med dette. Hva sier det om løsningene?
\begin{tasks}(3)
  \task $0 = 0$
  \task $0 = 7$
  \task $3x = 15$
\end{tasks}

\nivaahode{2}

\oppg{1.4}{2}{\uh}
Er likningen en identitet? Finn et moteksempel, eller vis med algebra at den alltid stemmer.
\begin{tasks}(1)
  \task $(x+5)^2 = x^2 + 25$
  \task $(x-4)(x+4) = x^2 - 16$
  \task $x^2 + 2 = 3x$. Den stemmer for både $x = 1$ og $x = 2$.
\end{tasks}

\oppg{1.5}{2}{\uh}
Løs likningen. Skriv om det er én løsning, ingen løsning eller alle tall.
\begin{tasks}(1)
  \task $2(3x-1) = 6x - 2$
  \task $5x - 2(x+1) = 3x + 4$
  \task $4(x-2) + x = 3x + 6$
\end{tasks}

\oppg{1.6}{2}{\uh}
Finn \textbf{den første feilen}. Løs riktig.
\[ 3(x+2) - (x-4) = 2x + 10 \ \to\ 3x + 6 - x + 4 = 2x + 10 \ \to\ 2x + 10 = 2x + 10 \ \to\ 0 = 0 \ \to\ x = 0 \]

\nivaahode{3}

\oppg{1.7}{3}{\uh}
Er likningen alltid sann (for alle tillatte verdier av $x$), noen ganger sann eller aldri sann? Begrunn svaret.
\begin{tasks}(1)
  \task $(x+6)^2 - (x-6)^2 = 24x$
  \task $x^2 = 2x$
  \task $x^2 + 1 = 0$
\end{tasks}

\oppg{1.8}{3}{\uh}
Lag tre likninger. Alle skal ha $x$ på begge sider og minst én brøk.
\begin{tasks}(1)
  \task En likning med nøyaktig én løsning, $x = 3$.
  \task En likning uten løsning.
  \task En identitet.
\end{tasks}
Kontroller hver likning ved å løse den.

\oppg{1.9}{3}{\uh}
Likningen $\sqrt{x^2} = x$ har uendelig mange løsninger. Er den sann for alle tillatte verdier av $x$, altså en identitet? Begrunn.

\begin{sjekk}
\begin{tasks}(2)
  \task Er $7x + 2$ et uttrykk eller en likning?
  \task Løs $5x + 4 = 5x - 1$.
  \task! Hvorfor kan vi aldri bevise en identitet ved å sette inn tall?
\end{tasks}
\end{sjekk}

% ---------------------------------------------------------------------

\del{2}{Likninger med brøk og formler}

\begin{ordliste}
fellesnevner & Et tall alle nevnerne går opp i. Vi bruker den minste. & \mittsprak\\
formel & En likning som viser sammenhengen mellom størrelser, for eksempel $A = g \cdot h$. & \mittsprak\\
variabel & En størrelse som kan ha ulike verdier, for eksempel tiden $t$. & \mittsprak\\
konstant & En størrelse som har samme verdi hele tiden, for eksempel prisen per kilo i $P = 40k$. & \mittsprak\\
snu en formel & Løse formelen med hensyn på en annen variabel. & \mittsprak\\
\end{ordliste}

\begin{mustkunne}
\begin{minipage}[t]{0.52\linewidth}
\textbf{Likninger med brøk.} Gang \emph{alle ledd} på begge sider med fellesnevneren. Da forsvinner brøkene.
\[ 6\cdot\frac{x+1}{2} = \frac{6}{1}\cdot\frac{x+1}{2} = 3(x+1) \]
Teller med flere ledd får parentes. Kontroller svaret i den \emph{opprinnelige} likningen.
\end{minipage}\hfill
\begin{minipage}[t]{0.44\linewidth}
\textbf{Snu en formel.}\par
\begin{enumerate}[leftmargin=1.4em,itemsep=1pt,topsep=2pt]
  \item Hva vil du ha alene?
  \item Gjør det samme med \emph{hele} venstre og høyre side.
  \item Fjern nevner først, deretter parentes.
  \item Kontroller med tall.
\end{enumerate}
\end{minipage}
\end{mustkunne}

\begin{eksempel}[1 — likning med brøk]
Løs $\dfrac{x-1}{3} - \dfrac{x+2}{4} = 1$.

\begin{stegtabell}
$12\cdot\dfrac{x-1}{3} - 12\cdot\dfrac{x+2}{4} = 12\cdot 1$ & Fellesnevner 12. Gang alle ledd.\\
$4(x-1) - 3(x+2) = 12$ & Parentes rundt tellerne.\\
$4x - 4 - 3x - 6 = 12$ & Gang ut. Minus gjelder begge ledd i parentesen.\\
$x - 10 = 12$, \ $x = 22$ & Slå sammen og løs.\\
Kontroll: & $\frac{21}{3} - \frac{24}{4} = 7 - 6 = 1$. Det stemmer.\\
\end{stegtabell}
\end{eksempel}

\begin{eksempel}[2 — snu en formel]
Trapesets areal er $A = \dfrac{(a+b)h}{2}$. Finn $b$ når $A = 30$, $a = 4$ og $h = 5$.

\begin{stegtabell}
$2A = (a+b)h$ & Gang begge sider med 2.\\
$\dfrac{2A}{h} = a + b$ & Del begge sider med $h$.\\
$b = \dfrac{2A}{h} - a$ & Trekk fra $a$ på begge sider.\\
$b = \dfrac{2\cdot 30}{5} - 4 = 8$ & Sett inn tallene.\\
Kontroll: & $\frac{(4+8)\cdot 5}{2} = 30$. Det stemmer.\\
\end{stegtabell}
\end{eksempel}

\begin{vanligfeil}
\begin{itemize}[leftmargin=1.2em,itemsep=1pt,topsep=0pt]
  \item $12\left(\frac{x-1}{3} - \frac{x+2}{4}\right) = 4(x-1) - \frac{x+2}{4}$ er feil. Alle ledd skal ganges: $4(x-1) - 3(x+2)$.
  \item $12\cdot\frac{x+2}{4} = 3x + 2$ er feil. Det mangler parentes: $3(x+2) = 3x + 6$.
  \item $b = \frac{2A}{h} + a$ er feil. Du må trekke fra $a$: $b = \frac{2A}{h} - a$.
\end{itemize}
\end{vanligfeil}

\nivaahode{1}

\oppg{2.1}{1}{\uh}
Løs likningen.
\begin{tasks}(3)
  \task $\dfrac{x}{3} + 2 = 5$
  \task $\dfrac{x}{2} = \dfrac{x}{5} + 6$
  \task $\dfrac{x+5}{3} = 4$
\end{tasks}

\oppg{2.2}{1}{\uh}
Vi skal løse $\dfrac{x+3}{2} + \dfrac{x}{3} = 4$.
\begin{tasks}(1)
  \task Fellesnevneren er 6. Fullfør: $6\cdot\dfrac{x+3}{2} + 6\cdot\dfrac{x}{3} = 6\cdot 4$ gir $3(x+3) + $\skrivlinje[2cm]$ = 24$.
  \task Løs likningen og kontroller svaret.
\end{tasks}

\oppg{2.3}{1}{\uh}
Formelen $F = 1{,}8C + 32$ gir temperaturen i fahrenheit når temperaturen i celsius er $C$.
\begin{tasks}(1)
  \task Hvilke størrelser er variable, og hvilke er konstante?
  \task Finn $F$ når $C = 25$.
  \task Finn $C$ når $F = 50$.
\end{tasks}

\oppg{2.4}{1}{\uh}
Farten er $v = \dfrac{s}{t}$, der $s$ er strekning og $t$ er tid.
\begin{tasks}(2)
  \task Løs formelen med hensyn på $s$.
  \task Løs formelen med hensyn på $t$.
\end{tasks}

\nivaahode{2}

\oppg{2.5}{2}{\uh}
Løs likningen. Kontroller svaret.
\begin{tasks}(1)
  \task $\dfrac{2x-1}{3} - \dfrac{x+4}{6} = 2$
  \task $\dfrac{x}{2} - \dfrac{x-3}{5} = 3$
  \task $\dfrac{x+3}{4} - \dfrac{x-2}{3} = \dfrac{1}{2}$
\end{tasks}

\oppg{2.6}{2}{\uh}
Finn \textbf{den første feilen}. Løs riktig.
\[ \frac{x+6}{2} - \frac{x-2}{3} = 4 \ \to\ 3(x+6) - 2(x-2) = 24 \ \to\ 3x + 18 - 2x - 4 = 24 \ \to\ x + 14 = 24 \ \to\ x = 10 \]

\oppg{2.7}{2}{\uh}
Løs formelen med hensyn på den variabelen som står i parentes.
\begin{tasks}(3)
  \task $O = 2(l + b)$. \ $(b)$
  \task $s = \dfrac{(v+u)t}{2}$. \ $(u)$
  \task $V = \dfrac{G \cdot h}{3}$. \ $(h)$
\end{tasks}

\oppg{2.8}{2}{\mh}
Det koster 15 kr å låse opp en elsparkesykkel. Deretter koster det 2 kr per minutt.
\begin{tasks}(1)
  \task Sett opp en formel for prisen $P$ kroner etter $t$ minutter. Hva er variabel, og hva er konstant?
  \task Tegn grafen i GeoGebra. Hvor mange minutter kan du kjøre for 95 kr?
  \task Løs formelen med hensyn på $t$, og kontroller svaret fra b).
\end{tasks}

\nivaahode{3}

\oppg{2.9}{3}{\uh}
Eli og Marte løser $\dfrac{x-4}{3} = \dfrac{x-4}{6}$.
Eli deler begge sider på $x - 4$ og får $\dfrac{1}{3} = \dfrac{1}{6}$. Hun svarer at det ikke er noen løsning.
Marte ganger med 6 og får en løsning. Hvem har rett? Forklar hva Eli glemte.

\oppg{2.10}{3}{\uh}
Se på formelen $\dfrac{x}{a} + \dfrac{x}{b} = 1$, der $a$ og $b$ er konstanter ulike fra 0.
\begin{tasks}(1)
  \task Løs formelen med hensyn på $x$.
  \task Hva skjer hvis $b = -a$? Forklar.
\end{tasks}
\hint{Når $x$ står i to ledd, kan du sette $x$ utenfor en parentes.}

\begin{sjekk}
\begin{tasks}(2)
  \task Løs $\dfrac{x}{4} + 1 = \dfrac{x}{2}$.
  \task Løs $F = m \cdot a$ med hensyn på $a$.
  \task! Hva blir $6\cdot\dfrac{x+7}{2}$? Forklar hvorfor det må være parentes.
\end{tasks}
\end{sjekk}
